\documentclass[12pt]{article}
\usepackage{amsmath, amssymb}
\usepackage{geometry}
\geometry{margin=1in}

\title{CSE400 -- Fundamentals of Probability in Computing\\
Lecture 5 Scribe: Bayes’ Theorem, Random Variables, and PMF}
\author{}
\date{}

\begin{document}

\maketitle

\section*{1. Bayes’ Theorem -- Weighted Average of Conditional Probabilities}

Let $A$ and $B$ be events. Event $A$ can be decomposed as
\[
A = AB \cup AB^c
\]
since any outcome in $A$ must either belong to both $A$ and $B$ or to $A$ but not to $B$.

Because $AB$ and $AB^c$ are mutually exclusive, by Axiom 3:
\[
Pr(A) = Pr(AB) + Pr(AB^c)
\]

Using conditional probability:
\[
Pr(AB) = Pr(A|B)Pr(B), \quad
Pr(AB^c) = Pr(A|B^c)Pr(B^c)
\]

Thus,
\[
Pr(A) = Pr(A|B)Pr(B) + Pr(A|B^c)[1 - Pr(B)]
\]

\textbf{Result:} $Pr(A)$ is a weighted average of conditional probabilities.

\section*{2. Learning by Example -- Insurance Model}

\subsection*{Example 3.1 (Part 1)}

Given:
\[
Pr(A_1|A)=0.4, \quad Pr(A_1|A^c)=0.2, \quad Pr(A)=0.3
\]

Let $A_1$ be the event of an accident within a year and $A$ the event of being accident prone.

\[
Pr(A_1)=Pr(A_1|A)Pr(A)+Pr(A_1|A^c)Pr(A^c)
\]
\[
= (0.4)(0.3) + (0.2)(0.7) = 0.26
\]

\subsection*{Example 3.1 (Part 2)}

\[
Pr(A|A_1)=\frac{Pr(A_1|A)Pr(A)}{Pr(A_1)}
= \frac{(0.4)(0.3)}{0.26} = \frac{6}{13}
\]

\section*{3. Law of Total Probability}

If $B_1, B_2, \dots, B_n$ form a partition of the sample space:
\[
Pr(A)=\sum_i Pr(A|B_i)Pr(B_i)
\]

\section*{4. Bayes’ Formula}

Using $Pr(AB_i)=Pr(B_i|A)Pr(A)$,
\[
Pr(B_i|A)=\frac{Pr(A|B_i)Pr(B_i)}{\sum_j Pr(A|B_j)Pr(B_j)}
\]

Where:
\[
Pr(B_i) = \text{a priori probability}, \quad
Pr(B_i|A) = \text{posterior probability}
\]

\section*{5. Example 3.2 -- Cards Problem}

Cards: RR, BB, RB.

Let $R$ be the event the upturned side is red.

\[
Pr(RB|R)=\frac{Pr(R|RB)Pr(RB)}{Pr(R|RR)Pr(RR)+Pr(R|RB)Pr(RB)+Pr(R|BB)Pr(BB)}
\]
\[
=\frac{(1/2)(1/3)}{(1)(1/3)+(1/2)(1/3)+(0)(1/3)}=\frac{1}{3}
\]

\section*{6. Random Variables}

A \textbf{random variable} is a real-valued function defined on the sample space. Probabilities are assigned to its values.

\subsection*{Example: Three Coin Tosses}

Let $Y$ = number of heads.

\[
P(Y=0)=\frac{1}{8}, \quad
P(Y=1)=\frac{3}{8}, \quad
P(Y=2)=\frac{3}{8}, \quad
P(Y=3)=\frac{1}{8}
\]
\[
\sum P(Y=y)=1
\]

\section*{7. Probability Mass Function (PMF)}

A random variable with countable values is discrete.

For discrete $X$ with range $R_X = \{x_1,x_2,\dots\}$:
\[
p(x_k)=Pr(X=x_k)
\]
\[
\sum_k p(x_k)=1
\]

\subsection*{PMF Example}

Given
\[
p(i)=c\lambda^i, \quad i=0,1,2,\dots
\]
\[
\sum_{i=0}^{\infty} c\lambda^i = 1
\]

Used to compute:
\[
P(X=0), \quad P(X>2)
\]

\end{document}

